% Generated draft: requires mathematical and source review.
\documentclass[11pt]{article}
\usepackage[margin=25mm]{geometry}
\usepackage{fontspec}
\setmainfont{Latin Modern Roman}
\usepackage{amsmath,amssymb,mathtools,unicode-math}
\setmathfont{Latin Modern Math}
\usepackage{enumitem,needspace,xurl,hyperref}
\setlength{\parindent}{0pt}
\setlength{\parskip}{5pt}
\setlength{\emergencystretch}{3em}
\newcommand{\N}{\mathbb{N}}
\newcommand{\Z}{\mathbb{Z}}
\newcommand{\Q}{\mathbb{Q}}
\newcommand{\R}{\mathbb{R}}
\newcommand{\C}{\mathbb{C}}
\newcommand{\F}{\mathbb{F}}
\newcommand{\A}{\mathbb{A}}
\newcommand{\PP}{\mathbb{P}}
\newcommand{\E}{\mathbb{E}}
\newcommand{\eps}{\varepsilon}
\newcommand{\dd}{\,\mathrm{d}}
\begin{document}
\section*{UnsolvedMath ID: 30000041}
\textbf{Source problem code:} catalogue entry
\subsection{Definitions and mathematical statement}
For $g\in L^1(\R)$ define convolution and the idempotency defect by
\[
 (g*g)(x)=\int_\R g(y)g(x-y)\,dy,
 \qquad J(g)=\|g*g-g\|_{L^1(\R)}.
\]
Young's inequality makes this a finite quantity. The source asks for its infimum among symmetric functions with $\|g\|_1=1$, and whether the infimum is attained.

Under the even-function reading, and for a specified scalar field $\mathbb K=\R$ or $\C$, this is the problem
\[
 c_{\mathbb K}=\inf\{J(g):g\in L^1(\R;\mathbb K),\
             g(-x)=g(x)\text{ a.e.},\ \|g\|_1=1\},
\]
and whether there exists an admissible $g_*$ with $J(g_*)=c_{\mathbb K}$. The source does not specify real versus complex values or whether symmetric instead means the involutive condition $g(-x)=\overline{g(x)}$. It also does not impose $g\geq0$. Unit $L^1$ norm alone is not the condition $\int g=1$ for signed or complex functions. A probability-density version would be a different explicitly restricted class and is not silently imposed.

The triage checked 21 August 2026 reports the original lower bound and a numerical Gaussian candidate, with no proof of the exact infimum or attainment located. Neither the bound's numerical value nor the candidate formula is supplied here. A positive even Gaussian gives only a test in any larger signed/complex admissible class, not by itself a global minimizer theorem. Constant spatial dilations $g(x)\mapsto a g(ax)$ preserve both $\|g\|_1$ and $J(g)$, so scale choices of such a candidate would also need a consistent interpretation. This does not settle existence of any minimizer.

The requested outputs remain an exact value, or sharp characterization, of the infimum and a separate proof of attainment or nonattainment in the intended class. The dated failure to locate a proof is recorded as supplied literature triage, not a new status verification.
\subsection{Short English statement}
The question minimizes a convolution defect on a unit $L^1$ sphere and separately asks about attainment. The intended symmetry and whether signed or complex functions are allowed must be fixed. The supplied lower-bound and numerical Gaussian information gives no proven sharp value or minimizer.
\subsection{Sources}
\begin{itemize}
\item Local UnsolvedMath registry, permanent ID 30000041, source code catalogue entry, supplied Oberwolfach Reports record; updated 16 September 2026.
\item \emph{Mini-Workshop: Wavelets and Frames (2004)}, Oberwolfach Reports, DOI 10.4171/owr/2004/10, pp. 479--544: \url{https://doi.org/10.4171/owr/2004/10}.
\item Supplied literature reference ``OWR 10/2004'': \url{https://doi.org/10.4171/OWR/2004/10}.
\item Supplied ``Feichtinger survey'': \url{https://doi.org/10.1007/s43670-024-00095-9}.
\item Literature triage checked 21 August 2026, derived from the supplied OWR curation; independently verify time-sensitive status and cited claims. No numerical bound or full candidate data are supplied.
\end{itemize}
\end{document}
